\documentclass{article}
\usepackage[landscape, margin=0.28cm]{geometry}
\usepackage{multicol}
\usepackage{amssymb}
\usepackage{amsmath}
\usepackage{enumitem}
\usepackage[compact]{titlesec}
\usepackage{lmodern}
\usepackage{mathpazo}
\renewcommand{\familydefault}{\sfdefault}
\usepackage{bm}
\usepackage{ctex}
\pagenumbering{gobble}

\titleformat{\section}{\normalfont\fontsize{9.2pt}{10pt}\bfseries}{\thesection}{0.18em}{}
\titleformat{\subsection}{\normalfont\fontsize{8.2pt}{9pt}\bfseries}{\thesubsection}{0.18em}{}
\titleformat{\subsubsection}{\normalfont\fontsize{7.4pt}{8pt}\bfseries}{\thesubsubsection}{0.18em}{}
\titlespacing*{\section}{0pt}{2pt}{1pt}
\titlespacing*{\subsection}{0pt}{1.5pt}{0.5pt}
\titlespacing*{\subsubsection}{0pt}{1pt}{0pt}
\setlength{\columnsep}{0.32cm}
\setlength{\columnseprule}{0.18pt}
\setlength{\parindent}{0pt}
\setlength{\parskip}{0pt}
\setlist[itemize]{nosep, leftmargin=*, label=\textbullet}
\setlist[enumerate]{nosep, leftmargin=*}
\allowdisplaybreaks
\setlength{\abovedisplayskip}{0.5pt}
\setlength{\belowdisplayskip}{0.5pt}

\newcommand{\key}[1]{\textbf{#1}}
\newcommand{\lowterm}[1]{\textbf{\underline{#1}}}
\newcommand{\term}[1]{\textit{#1}}
\newcommand{\R}{\mathbb{R}}
\newcommand{\E}{\mathbb{E}}
\newcommand{\N}{\mathcal{N}}
\newcommand{\KL}{\mathrm{KL}}
\newcommand{\1}{\mathbf 1}
\DeclareMathOperator*{\argmax}{arg\,max}
\DeclareMathOperator*{\argmin}{arg\,min}
\DeclareMathOperator{\softmax}{softmax}
\DeclareMathOperator{\sigmoid}{sigmoid}
\DeclareMathOperator{\tr}{tr}

\begin{document}
\tiny
\begin{multicols*}{4}

\section{L2--3 Supervised Learning}
\subsection{Classification and Loss}
\key{ERM}: population risk $\theta^\star=\arg\min_\theta\E_{x\sim P}\operatorname{err}(f(x;\theta),y(x))$; empirical risk $\hat\theta=\arg\min_\theta\frac1N\sum_i\operatorname{err}(f(x_i;\theta),y_i)$. \key{Binary logistic}: $p(y=1|x)=\sigma(w^\top x+b)$,
\[
\ell=-y\log p-(1-y)\log(1-p),\quad \frac{\partial\ell}{\partial z}=p-y.
\]
\key{Multi-class}: logits $z\in\R^K$, $p_k=e^{z_k}/\sum_j e^{z_j}$,
\[
\ell(z,y)=-\log p_y=-z_y+\log\sum_j e^{z_j},\quad
\frac{\partial\ell}{\partial z_k}=p_k-\1[k=y].
\]
\lowterm{One-hot CE}: $-\sum_k y_k\log p_k$.

\subsection{MLP and Backprop}
\key{Layer}: $h^{(0)}=x$, $z^{(\ell)}=W^{(\ell)}h^{(\ell-1)}+b^{(\ell)}$, $h^{(\ell)}=\phi(z^{(\ell)})$. Batch code convention often uses $X\in\R^{B\times n}$, $Z=XW+b$, $W\in\R^{n\times m}$.
\key{Backprop}: define $\delta^{(\ell)}=\partial L/\partial z^{(\ell)}$:
\[
\partial_{W^{(\ell)}}L=\delta^{(\ell)}(h^{(\ell-1)})^\top,\quad
\partial_{b^{(\ell)}}L=\delta^{(\ell)},
\]
\[
\delta^{(\ell-1)}=((W^{(\ell)})^\top\delta^{(\ell)})\odot\phi'(z^{(\ell-1)}).
\]
\key{Activation}: sigmoid/tanh saturate; ReLU $\max(0,z)$ keeps positive-side gradient; ELU/smooth variants improve negative side. \lowterm{Subgradient}: ReLU at $0$ uses any value in $[0,1]$ by convention.

\subsection{CNN}
\key{Motivation}: locality + weight sharing. Convolution:
\[
y_{i,j,o}=\sum_{u,v,c}K_{u,v,c,o}x_{i+u,j+v,c}+b_o.
\]
For input $H\times W\times C_{\rm in}$, kernel $K_h\times K_w$, padding $P$, stride $S$:
\[
H'=\lfloor(H+2P-K_h)/S\rfloor+1,\quad
W'=\lfloor(W+2P-K_w)/S\rfloor+1.
\]
Params $K_hK_wC_{\rm in}C_{\rm out}+C_{\rm out}$. \lowterm{Scanning MLP}: same detector slides across space. \lowterm{Pooling}: max/avg reduces spatial size and jitter. \lowterm{Padding}: keeps boundary info; same padding for odd kernel uses $P=(K-1)/2$.

\subsection{Optimization}
\key{GD}: $x_{k+1}=x_k-\eta_k\nabla f(x_k)$. If $\nabla f$ is $L$-Lipschitz:
\[
f(y)\le f(x)+\nabla f(x)^\top(y-x)+\frac L2\|y-x\|^2,
\]
so $0<\eta<2/L$ decreases the upper bound. \key{Newton}: from quadratic Taylor,
\[
x_{k+1}=x_k-\eta[\nabla^2f(x_k)]^{-1}\nabla f(x_k).
\]
\key{Momentum}: $v_{k+1}=\beta v_k+g_k$, $x_{k+1}=x_k-\eta v_{k+1}$. \key{Nesterov}: compute gradient at lookahead $x_k+\beta(x_k-x_{k-1})$.
\key{AdaGrad}: $s_k=s_{k-1}+g_k^2$, $x_{k+1}=x_k-\eta g_k/(\sqrt{s_k}+\epsilon)$; may decay too fast. \key{RMSProp}: $s_k=\rho s_{k-1}+(1-\rho)g_k^2$. \key{Adam}:
\[
m_k=\beta_1m_{k-1}+(1-\beta_1)g_k,\quad v_k=\beta_2v_{k-1}+(1-\beta_2)g_k^2,
\]
\[
x_{k+1}=x_k-\eta\hat m_k/(\sqrt{\hat v_k}+\epsilon),\quad
\hat m_k=m_k/(1-\beta_1^k).
\]

\subsection{Regularization and Architectures}
\lowterm{Dropout}: randomly zero units, approximates ensemble and reduces co-adaptation. \lowterm{Weight decay}: add $\lambda\|\theta\|^2$. \lowterm{Gradient clipping}: $g\leftarrow g\min(1,\tau/\|g\|)$.
\key{BatchNorm}: $\mu_B=\frac1B\sum_i z_i$, $\sigma_B^2=\frac1B\sum_i(z_i-\mu_B)^2$,
\[
\hat z_i=\gamma\frac{z_i-\mu_B}{\sqrt{\sigma_B^2+\epsilon}}+\beta.
\]
\key{LayerNorm}: normalize hidden dimension inside one sample; better for sequence models. \key{ResNet}: $h_{\ell+1}=h_\ell+F(h_\ell)$, residual path eases optimization. \key{DenseNet}: concatenate previous features. \key{FCN}: replace fully connected layer by convolution for dense prediction.

\section{L4 Energy-Based Models}
\subsection{Energy View}
\key{Generative modeling}: learn $P(X)$ or $P(X,y)=P(y)P(X|y)$ instead of $P(y|X)$. \lowterm{EBM}: unnormalized score via energy
\[
p_\theta(x)=\frac{e^{-E_\theta(x)}}{Z_\theta},\quad Z_\theta=\int e^{-E_\theta(x)}dx.
\]
Low energy means high probability; training and inference both need sampling or partition-function gradients.

\subsection{Hopfield Network}
State $y_i\in\{-1,1\}$, local field $h_i=\sum_{j\ne i}w_{ij}y_j+b_i$. Update flips if $y_ih_i<0$.
\[
E(y)=-\sum_{i<j}w_{ij}y_iy_j-\sum_i b_iy_i=-\frac12y^\top Wy-b^\top y.
\]
Asynchronous flip changes energy by $\Delta E=2y_ih_i$; if $y_ih_i<0$, energy decreases. Finite states imply convergence to a local minimum. \key{Hebbian storage}: $w_{ij}=\frac1P\sum_p y_i^py_j^p$, $w_{ii}=0$. Multiple patterns cause \lowterm{spurious minima}; capacity is limited, roughly linear in neuron count for stable random memories.
\key{Optimization view}: desired patterns should be valleys; short-run evolution can raise undesired valleys by contrastive updates.

\subsection{Boltzmann Machine}
\key{Stochastic Hopfield}: $p(y)\propto e^{-E(y)}$. Conditional Gibbs update for $\{-1,1\}$ units:
\[
P(y_i=1|y_{\neg i})=\sigma(2h_i).
\]
Temperature $T$: $p_T(y)\propto e^{-E(y)/T}$; annealing lowers $T$ from exploration to exploitation. \key{MLE gradient}:
\[
\nabla_{w_{ij}}\log p_\theta(y)=y_iy_j-\E_{p_\theta}[y_iy_j].
\]
Learning = \lowterm{positive phase} data correlation minus \lowterm{negative phase} model correlation. With hidden $h$:
\[
\nabla L=\E_{p(h|v_{\rm data})}[y_iy_j]-\E_{p(v,h)}[y_iy_j].
\]

\subsection{RBM and Contrastive Divergence}
\lowterm{RBM}: bipartite $v,h$; no visible-visible or hidden-hidden edges:
\[
E(v,h)=-v^\top Wh-b^\top v-c^\top h.
\]
Conditionals factorize:
\[
P(h_j=1|v)=\sigma(c_j+W_{:j}^\top v),\quad
P(v_i=1|h)=\sigma(b_i+W_{i:}h).
\]
\key{CD-$k$}: start at data $v^{(0)}$, run $k$ Gibbs steps, update $W\propto v^{(0)}h^{(0)\top}-v^{(k)}h^{(k)\top}$.

\subsection{Sampling}
\key{Inverse CDF}: sample $u\sim U[0,1]$, choose category by cumulative mass. \key{Box-Muller}: Gaussian from uniform:
\[
z_1=\sqrt{-2\log u_1}\cos(2\pi u_2),\quad z_2=\sqrt{-2\log u_1}\sin(2\pi u_2).
\]
\key{Importance sampling}: $\E_p[f(x)]=\E_q[f(x)p(x)/q(x)]$. Good $q$ must cover $p$'s high-mass region.
\key{MCMC}: Markov chain with stationary $\pi$. Detailed balance $\pi(x)T(x\to x')=\pi(x')T(x'\to x)$ implies stationarity. \key{MH}: proposal $q(x'|x)$, accept
\[
\alpha=\min\left(1,\frac{\pi(x')q(x|x')}{\pi(x)q(x'|x)}\right).
\]
For EBM, partition $Z$ cancels. \key{Gibbs}: update one conditional exactly. \key{SGLD}: $x_{t+1}=x_t-\frac{\eta}{2}\nabla_xE_\theta(x_t)+\sqrt{\eta}\xi_t$.

\section{L5 Generative Adversarial Networks}
\subsection{Implicit Generative Model}
Instead of density $p_\theta(x)$, learn sampler $x=G_\theta(z)$, $z\sim p_z$. Need differentiable distance between generated and real distributions; GAN learns this distance with a discriminator.
\[
\min_G\max_D V(D,G)=\E_{x\sim p_{\rm data}}\log D(x)+\E_{z\sim p_z}\log(1-D(G(z))).
\]
For fixed $G$:
\[
D^\star(x)=\frac{p_{\rm data}(x)}{p_{\rm data}(x)+p_g(x)}.
\]
Substitute:
\[
V(D^\star,G)=-\log4+2\,{\rm JSD}(p_{\rm data}\|p_g).
\]
Non-saturating generator loss $-\E_z\log D(G(z))$ gives stronger gradients.

\subsection{Evaluation and Failure Modes}
\lowterm{Mode collapse}: generator covers few modes. \lowterm{Training instability}: moving target game; discriminator too strong $\Rightarrow$ weak generator gradient. \lowterm{IS}: rewards confident/diverse ImageNet predictions, but blind to intra-class artifacts. \lowterm{FID}: compare Gaussian feature statistics:
\[
\|\mu_r-\mu_g\|^2+\tr(\Sigma_r+\Sigma_g-2(\Sigma_r\Sigma_g)^{1/2}).
\]

\subsection{WGAN Family}
JSD can be flat when supports do not overlap. Wasserstein-1:
\[
W(p,q)=\sup_{\|f\|_L\le1}\E_{x\sim p}f(x)-\E_{x\sim q}f(x).
\]
WGAN critic maximizes $\E_{p_{\rm data}}D(x)-\E_zD(G(z))$. \lowterm{Lipschitz constraint}: weight clipping, spectral norm, or WGAN-GP
\[
\lambda\E_{\hat x}(\|\nabla_{\hat x}D(\hat x)\|_2-1)^2.
\]
\key{DCGAN}: conv/deconv, batchnorm, stable architecture. \key{BigGAN}: scale + class conditioning + truncation trick. \key{GigaGAN}: text-to-image GAN scaling. \key{R3GAN}: regularized robust objective.

\subsection{Variants and Applications}
\key{Conditional GAN}: $G(z,c)$, $D(x,c)$. \key{Semi-supervised GAN}: discriminator has $K+1$ classes, $K$ real labels plus fake. \key{BiGAN/ALI}: learn encoder $E(x)$ with discriminator over $(x,z)$ pairs; gives representation learning. \key{Pix2Pix}: paired image-to-image translation with conditional GAN + reconstruction loss. \key{CycleGAN}: unpaired translation using cycle consistency
\[
G:X\to Y,\quad F:Y\to X,\quad \|F(G(x))-x\|+\|G(F(y))-y\|.
\]
\key{DragGAN}: interactive image manipulation by optimizing latent/code with handle-point constraints.

\section{L6 Flow and VAE}
\subsection{Normalizing Flow}
Invertible differentiable $x=G_\theta(z)$:
\[
\log p_X(x)=\log p_Z(G^{-1}(x))+\log\left|\det\frac{\partial G^{-1}}{\partial x}\right|.
\]
Composition $G=f_K\circ\cdots\circ f_1$ adds log determinants. Flow design needs invertible layer, efficient inverse, efficient determinant.
\key{Coupling}: split $x=(x_a,x_b)$:
\[
y_a=x_a,\quad y_b=x_b\odot e^{s(x_a)}+t(x_a),\quad
\log|\det J|=\sum s(x_a).
\]
\key{Glow}: invertible $1\times1$ convolution mixes channels; LU decomposition makes determinant cheap: $\log|\det W|=\sum_i\log|u_{ii}|$. \key{Dequantization}: add uniform noise to discrete pixels so continuous density is well-defined.
\key{Continuous NF / Neural ODE}: $\frac{dz(t)}{dt}=f_\theta(z(t),t)$,
\[
\log p(z(t_1))=\log p(z(t_0))-\int_{t_0}^{t_1}\tr\left(\frac{\partial f_\theta}{\partial z(t)}\right)dt.
\]

\subsection{EM and Variational Inference}
Latent-variable likelihood $\log p_\theta(x)=\log\sum_z p_\theta(x,z)$. Jensen:
\[
\log p_\theta(x)=\log\E_q\frac{p_\theta(x,z)}{q(z)}
\ge \E_q\log p_\theta(x,z)+H(q)=\mathcal L(q,\theta).
\]
Gap:
\[
\log p_\theta(x)-\mathcal L(q,\theta)=\KL(q(z)\|p_\theta(z|x)).
\]
EM: E-step $q(z)=p_{\theta^{old}}(z|x)$, M-step maximize $\E_q\log p_\theta(x,z)$.

\subsection{VAE}
Encoder $q_\phi(z|x)$, decoder $p_\theta(x|z)$, prior $p(z)$. ELBO:
\[
\log p_\theta(x)\ge \E_{q_\phi(z|x)}[\log p_\theta(x|z)]-\KL(q_\phi(z|x)\|p(z)).
\]
Reparameterization:
\[
z=\mu_\phi(x)+\sigma_\phi(x)\odot\epsilon,\quad \epsilon\sim\N(0,I).
\]
Diagonal Gaussian KL:
\[
\KL(\N(\mu,\sigma^2I)\|\N(0,I))
=\frac12\sum_i(\mu_i^2+\sigma_i^2-\log\sigma_i^2-1).
\]
\key{Terms}: reconstruction = data fidelity; KL = latent regularization. \key{PixelVAE/NVAE}: hierarchical latent variables improve global coherence and likelihood.

\section{L7 Diffusion and Flow Matching}
\subsection{DDPM}
Forward:
\[
q(x_t|x_{t-1})=\N(\sqrt{\alpha_t}x_{t-1},\beta_tI),\quad \bar\alpha_t=\prod_{s=1}^t\alpha_s.
\]
Closed form:
\[
q(x_t|x_0)=\N(\sqrt{\bar\alpha_t}x_0,(1-\bar\alpha_t)I),\quad
x_t=\sqrt{\bar\alpha_t}x_0+\sqrt{1-\bar\alpha_t}\epsilon.
\]
Posterior $q(x_{t-1}|x_t,x_0)=\N(\tilde\mu_t,\tilde\beta_tI)$:
\[
\tilde\mu_t=\frac{\sqrt{\bar\alpha_{t-1}}\beta_t}{1-\bar\alpha_t}x_0+
\frac{\sqrt{\alpha_t}(1-\bar\alpha_{t-1})}{1-\bar\alpha_t}x_t,\quad
\tilde\beta_t=\frac{1-\bar\alpha_{t-1}}{1-\bar\alpha_t}\beta_t.
\]
Predict noise:
\[
x_0=\frac{x_t-\sqrt{1-\bar\alpha_t}\epsilon}{\sqrt{\bar\alpha_t}},\quad
L_{\rm simple}=\E\|\epsilon-\epsilon_\theta(x_t,t)\|^2.
\]
Sampling:
\[
x_{t-1}=\frac{1}{\sqrt{\alpha_t}}\left(x_t-\frac{\beta_t}{\sqrt{1-\bar\alpha_t}}\epsilon_\theta(x_t,t)\right)+\sigma_tz.
\]

\subsection{Score-Based Models}
\lowterm{Score}: $s_t(x)=\nabla_x\log p_t(x)$. Denoising score matching avoids unknown data score by using noisy conditional score:
\[
\nabla_{x_t}\log q(x_t|x_0)=-\frac{x_t-\sqrt{\bar\alpha_t}x_0}{1-\bar\alpha_t}.
\]
SDE:
\[
dx=f(x,t)dt+g(t)dW_t,\quad
dx=[f(x,t)-g(t)^2\nabla_x\log p_t(x)]dt+g(t)d\bar W_t.
\]
Probability-flow ODE:
\[
dx=\left[f(x,t)-\frac12g(t)^2\nabla_x\log p_t(x)\right]dt.
\]
Likelihood along ODE:
\[
\log p_0(x_0)=\log p_T(x_T)+\int_0^T\tr\left(\frac{\partial f_\theta}{\partial x_t}\right)dt.
\]
\key{Annealed Langevin}: use multiple noise levels; high noise gives global movement, low noise refines detail.

\subsection{DPM-Solver and Applications}
DPM-Solver integrates diffusion ODE with high-order exponential integrators, reducing sampling steps. \lowterm{Latent diffusion}: diffuse in autoencoder latent space, cheaper than pixel space. \key{DiT/SORA}: transformer over spacetime latent patches. \key{DALL-E 2}: CLIP prior + diffusion decoder. \key{DALL-E 3}: synthetic captions improve prompt following. \key{SDEdit}: add noise to input image then denoise under prompt. \key{ControlNet}: extra condition branch injects edges/depth/pose. \key{Logit-normal sampling}: biases timestep sampling toward informative middle region.

\subsection{Flow Matching}
Continuous flow $dx_t/dt=v_\theta(x_t,t)$. Objective:
\[
\mathcal L_{\rm FM}=\E_{t,x_t}\|v_\theta(x_t,t)-u_t(x_t)\|^2.
\]
Conditional path from noise $x_0$ to data $x_1$: $x_t=(1-t)x_0+tx_1$, target $u_t=x_1-x_0$. Conditional Flow Matching (CFM) samples pairs and regresses conditional velocity; marginal vector field is recovered in expectation. \key{Non-uniqueness}: many vector fields induce same marginals; choice affects training variance and sampling trajectory. \key{SD3}: flow target often parameterizes velocity between noise/data under continuous timestep.

\section{L8--9 Sequence Modeling}
\subsection{Sequence Data and Tokenization}
Sequence data includes language, audio, action; images/DNA/protein can be serialized. Autoregressive factorization:
\[
p(x_{1:T})=\prod_{t=1}^T p(x_t|x_{<t}).
\]
\key{Tokenization}: text $\to$ basic units $\to$ token IDs. Byte-Pair Encoding (BPE) repeatedly merges frequent adjacent symbols; balances vocabulary size and sequence length. \key{Embedding}: token ID $i\to e_i=W_E[i]$; positional info is needed because token embeddings alone are order-free.

\subsection{Autoregressive CNN and Pixel Models}
\key{WaveNet}: causal/dilated temporal convolution expands receptive field without recurrence. \key{PixelCNN}: image likelihood
\[
p(x)=\prod_i p(x_i|x_{<i})
\]
under raster order; masked convolution prevents access to future pixels. Gated PixelCNN uses gates to improve expressiveness.

\subsection{RNN and LSTM}
\key{RNN}:
\[
h_t=\tanh(W_{xh}x_t+W_{hh}h_{t-1}+b_h),\quad y_t=W_{hy}h_t+b_y.
\]
BPTT multiplies many Jacobians; singular values $<1$ vanish, $>1$ explode. \key{LSTM}:
\[
f_t=\sigma(W_{xf}x_t+W_{hf}h_{t-1}+b_f),\quad
i_t=\sigma(W_{xi}x_t+W_{hi}h_{t-1}+b_i),
\]
\[
\tilde c_t=\tanh(W_{xc}x_t+W_{hc}h_{t-1}+b_c),\quad
c_t=f_t\odot c_{t-1}+i_t\odot\tilde c_t,
\]
\[
o_t=\sigma(W_{xo}x_t+W_{ho}h_{t-1}+b_o),\quad
h_t=o_t\odot\tanh(c_t).
\]
Additive cell path gives controllable memory; forget gate decides retention.

\subsection{Attention and Transformer}
For token $t$: $q_t=W_Qx_t$, $k_j=W_Kx_j$, $v_j=W_Vx_j$,
\[
\alpha_{tj}=\frac{\exp(q_t^\top k_j/\sqrt{d_k})}{\sum_i\exp(q_t^\top k_i/\sqrt{d_k})},\quad
y_t=\sum_j\alpha_{tj}v_j.
\]
Matrix form:
\[
\operatorname{Attn}(Q,K,V)=\softmax(QK^\top/\sqrt{d_k})V.
\]
\key{Multi-head}: parallel attention subspaces, concatenate and project. \key{Causal mask}: block future tokens. \key{Transformer block}: attention + MLP + residual + normalization. Complexity $O(T^2d)$; FlashAttention reduces memory IO, Mamba/SSM explores linear-time sequence mixing.

\subsection{Objectives and Case Studies}
\key{Objectives}: autoregressive LM, masked/denoising, contrastive predictive. \key{Recipe}: pretraining, fine-tuning, low-rank tuning. \key{GPT}: decoder-only causal Transformer; next-token pretraining gives in-context learning with scale. \key{ViT}: image patches as tokens. \key{Swin}: shifted local windows for hierarchical vision. \key{CLIP}: contrastive image-text alignment,
\[
L=-\frac12\left[\frac1N\sum_i\log\frac{e^{s_{ii}/\tau}}{\sum_j e^{s_{ij}/\tau}}+
\frac1N\sum_i\log\frac{e^{s_{ii}/\tau}}{\sum_j e^{s_{ji}/\tau}}\right].
\]
\key{SigLIP}: pairwise sigmoid loss avoids global softmax. \key{Whisper}: audio Transformer trained on large-scale speech. \key{MAR}: masked autoregressive modeling for images.

\section{L10 Modern Large Models}
\subsection{Pretraining and Scaling}
\key{Autoregressive LM}: maximize $\sum_t\log p_\theta(x_t|x_{<t})$. Base model learns continuation; chat model learns instruction-following behavior through post-training. \key{Scaling laws}: loss decreases smoothly with model size $N$, data $D$, compute $C$, often power-law:
\[
L(N)\approx L_\infty+aN^{-\alpha},\quad
L(D)\approx L_\infty+bD^{-\beta}.
\]
Compute-optimal training balances parameters and tokens; not simply largest model on fixed data. \lowterm{Prompting}: zero-shot/few-shot describes task in context; no parameter update. \lowterm{Instruction tuning}: supervised fine-tuning on instruction-response mixtures.

\subsection{RLHF and Hallucination}
\key{SFT}: next-token imitation on human demonstrations. It cannot express preference among multiple valid answers. \key{Reward model}: train from pairwise rankings:
\[
L_{\rm RM}=-\log\sigma(r_\phi(x,y_w)-r_\phi(x,y_l)).
\]
\key{Policy optimization with KL}:
\[
\max_\pi\E_{y\sim\pi(\cdot|x)}[r_\phi(x,y)]-\beta\KL(\pi(\cdot|x)\|\pi_{\rm ref}(\cdot|x)).
\]
KL prevents policy from drifting too far from SFT/reference. RLHF can reward uncertainty awareness, e.g. ``I don't know'' when question is unanswerable. Negative samples provide learning signal unavailable in pure SFT.

\subsection{RAG and Reasoning}
\key{RAG}: retrieve documents $d$ then generate:
\[
p(y|x)\approx\sum_{d\in\mathcal R(x)}p_\theta(y|x,d)p(d|x).
\]
Retrieval supplies missing knowledge but does not automatically solve multi-step reasoning. \lowterm{Inference-time scaling}: spend more test-time compute on search, verification, self-consistency, or long reasoning trajectories. For reasoning model:
\[
\max_\pi \E_{\tau\sim\pi(\cdot|x)}[R(\operatorname{answer}(\tau),y)]-\beta\KL(\pi\|\pi_{\rm ref}).
\]
Answer verifier gives sparse reward; RL can optimize thinking trajectory even without unique chain-of-thought labels.

\subsection{Multimodal Models}
\key{Vision-language few-shot}: convert image/video tasks to text generation conditioned on visual tokens. \key{Flamingo}: frozen vision encoder + frozen LM with cross-attention layers and Perceiver resampler for interleaved image/video-text. \key{LLaVA}: CLIP vision encoder + projection to Vicuna/LLM; visual instruction tuning teaches chat-style responses. \key{Unified multimodal model}: one decoder handles understanding, generation, editing, reasoning; may use separate experts/tokenizers for text vs image. \key{BAGEL-style idea}: decoder-only unified model with modality-specific experts; understanding often emerges before reasoning/editing, which need larger scale.

\section{Cross-Lecture Map}
\key{Density models}: EBM defines $p\propto e^{-E}$ but sampling is hard; flow gives exact likelihood via invertible map; VAE gives ELBO; diffusion learns reverse denoising; flow matching learns vector field directly. \key{Implicit models}: GAN learns sampler without tractable likelihood. \key{Sequence models}: Transformer turns all modalities into token sequences. \key{Large models}: pretraining builds capability, post-training aligns behavior, inference-time compute and multimodal grounding extend usability.

\end{multicols*}
\end{document}
